\section{Risk-Controlled Context Selection}

\subsection{Setting}

Let $S_0$ be an anchor, $x$ a query, $y$ its target, and
$\candidates=\{C_1,\ldots,C_K\}$ a candidate pool of auxiliary contexts. A
predictor $f_A$ is trained on the anchor and the contexts indexed by
$A\subseteq\{1,\ldots,K\}$ and is kept fixed during selection. With $\loss$ the
end-task loss, the utility of a set and the marginal utility of one candidate
are
\begin{equation}
 F(A)=\E\!\left[\loss(f_{\emptyset}(x),y)-\loss(f_A(x),y)\right],
 \qquad
 \marginal(j\mid A)=F(A\cup\{j\})-F(A).
 \label{eq:marginal}
\end{equation}
A routing policy observes a state $A$ and a candidate score
$\widehat m(j\mid A)$ produced by a utility model, and returns either a
candidate to add or a decision to abstain. We write $z$ for a decision state,
which collects the current set, the candidate pool, the scores, and the
response features the model uses.

\subsection{Harm and its budget}

Routing is harmful when the added context increases the loss. Since the loss
is only observed after the fact, the harmful event is defined through the true
marginal,
\begin{equation}
 \harm(z,j)=\mathbf{1}\!\left\{\marginal(j\mid A)<0\right\}.
 \label{eq:harm}
\end{equation}
A decision rule that maps a state to a candidate $j(z)$ and a binary action
$a(z)\in\{0,1\}$ induces the harmful-selection mass
\begin{equation}
 R=\E\!\left[a(z)\,\harm\!\left(z,j(z)\right)\right],
 \label{eq:mass}
\end{equation}
that is, the expected number of harmful selections per decision. The rule's
coverage is $C=\Pr[a(z)=1]$, and among the decisions it accepts the harmful
share is
\begin{equation}
 \mathrm{SR}=\Pr\!\left[\harm(z,j(z))=1 \mid a(z)=1\right]=\frac{R}{C}.
 \label{eq:selective-risk}
\end{equation}
The router's objective is the constrained problem
\begin{equation}
 \max_{a,j}\ C
 \quad\text{subject to}\quad
 R\le\alpha ,
 \label{eq:objective}
\end{equation}
for a harm budget $\alpha$ chosen by the operator. Equation~\eqref{eq:objective}
separates two questions that a ranking model conflates: which candidate is
best, and whether any candidate should be used. A rule that always acts
maximises $C$ and says nothing about $R$; a rule that never acts has $R=0$ and
$C=0$. The interesting decisions lie between them.

A sequential router applies such a rule repeatedly. At each step it either
adds the selected candidate, abstains at that step but continues, or stops and
returns the current set. Stopping is the special case in which no remaining
candidate can be certified, and abstention at the first step corresponds to
using the anchor alone.
